% !TeX root = ../main.tex

\chapter{}



\section{补充章节}

附录对论文来讲主要起到补充说明的作用，附属于正文，必要时才添加，一般情况是省略掉的，
把原文中要说明的相关材料，最好是整理成表格放到论文的最后，并且标上标号，要把它在原文出现的地方找到，
在原文中写“见附录X”，这样读者就可以根据需要去查阅附录内容。附录的内容一般包括以下几类：

1、部分材料编入正文中会让文章显得主次不清，缺乏逻辑性，省略掉又会让文章显得不完整，
这类材料主要是一些比正文更为详细的研究方法和技术，对全文起到重要的补充作用。

2、部分资料由于篇幅过长，或者是复制品，不便于在正文展示，这时候我们就可以考虑使用附录进行补充说明。

3、一些对正文非常重要的原始数据、推导公式、源码程序、框图、统计表、设计图纸、调查问卷等，
这些部分不便于省略，我们就在正文后以附录的形式体现。

